% Extracting Steering Vectors from the J Space
% Reconstructed from the submitted PDF (102_Extracting_Steering_Vector.pdf).
% Compile with XeTeX via tectonic:  tectonic main.tex
\documentclass{article}

% During submission, before acceptance
\usepackage[dblblindworkshop]{neurips_2026}
% After acceptance, add "final" to compile the camera-ready version:
% \usepackage[dblblindworkshop, final]{neurips_2026}

% ---------------------------------------------------------------- fonts (XeTeX)
\usepackage{fontspec}
\setmainfont{Times New Roman}
% Chinese tokens (the submission used STSongti-SC-Regular = "Songti SC",
% scaled to 0.92: measured from the glyph boxes in the submitted PDF)
\newfontfamily\zhfont{Songti SC}[Scale=0.92]
\newcommand{\zh}[1]{{\zhfont #1}}

% ---------------------------------------------------------------- packages
\usepackage{hyperref}       % hyperlinks
\usepackage{url}            % simple URL typesetting
\usepackage{booktabs}       % professional-quality tables
\usepackage{amsmath}
\usepackage{amsfonts}       % blackboard math symbols
\usepackage{nicefrac}       % compact symbols for 1/2, etc.
\usepackage{microtype}      % microtypography
\usepackage[table]{xcolor}  % colors (+ colortbl)
\usepackage{array}
\usepackage{multirow}
\usepackage{graphicx}
\usepackage{listings}
\usepackage{caption}

% ---------------------------------------------------------------- colours
% table shading: green = ours ahead by > 0.10, red = fitted ahead by > 0.10
\definecolor{posgreen}{rgb}{0.88,1,0.88}   % #e0ffe0
\definecolor{negred}{rgb}{1,0.88,0.88}     % #ffe0e0
\definecolor{ctrlgrey}{gray}{0.5}          % #808080, J=I and random columns
% listings
\definecolor{lstbg}{gray}{0.98}
\definecolor{lstframe}{gray}{0.8}
\definecolor{lstkw}{rgb}{0,0,0.8}
\definecolor{lststr}{rgb}{0.7,0,0}
\definecolor{lstcmt}{rgb}{0,0.5,0}

\newcommand{\pos}{\cellcolor{posgreen}}
\newcommand{\negc}{\cellcolor{negred}}
\newcommand{\dlt}[1]{{\scriptsize(#1)}}
\newcolumntype{g}{>{\color{ctrlgrey}}r}

% Emoji placeholder. The submitted PDF prints the literal text "[emoji]" for
% every emoji concept token in Table 3; the actual emoji are not recoverable
% from the PDF.  TODO: replace with the real emoji tokens.
% In the submission the placeholder is set at 8pt (\footnotesize), larger than the
% surrounding 7pt token list; reproduced here.
\newcommand{\emojiph}[1][]{{\footnotesize\texttt{[emoji]}}}

\lstset{
  language=Python,
  basicstyle=\ttfamily\scriptsize,
  keywordstyle=\color{lstkw}\bfseries,
  keywordstyle={[2]\color{lstkw}\bfseries},
  stringstyle=\color{lststr},
  commentstyle=\color{lstcmt}\itshape,
  backgroundcolor=\color{lstbg},
  frame=single,
  rulecolor=\color{lstframe},
  xleftmargin=3pt,
  framexleftmargin=3pt,
  belowskip=\smallskipamount,
  showstringspaces=false,
  captionpos=t,
  morekeywords={True,False},
}

\title{Extracting Steering Vectors from the J Space}
\workshoptitle{Neural Network Artifacts as a New Data Modality}

\author{%
  Anonymous Author(s) \\
  Affiliation \\
  Address \\
  \texttt{email} \\
}

\begin{document}

\maketitle

\begin{abstract}
  In this work, we explore the idea of deriving steering vectors from a published Jacobian lens rather than from data. To that end, we run the lens backwards: the unembedding rows of a few concept tokens are pulled back through the lens Jacobian to the activation direction that raises their lens logit, contrasted against counter-tokens, and averaged. We evaluate the resulting direction against a difference-in-means vector fitted on the model over 8 behaviours on Qwen3-1.7B and show that it matches or exceeds the fitted vector on behaviours that amount to a preference over word choice. Finally, we present results suggesting that the method fails where a behaviour requires replacing the lexical stream: the derived direction encodes the topic of its concept tokens rather than the behaviour, and a projection of the fitted direction onto the subspace the tokens can reach fails on refusal and pirate register as well.
\end{abstract}

% =====================================================================
\section{Introduction}

Steering vectors are directions added to the residual stream of a language model that shift its behaviour, for instance towards answering in capitals or towards refusing. They are usually obtained by contrasting activations on paired prompts and taking a difference of means \citep{turner2023actadd,rimsky2024caa,zou2023repe}, which requires a dataset, a forward pass per example, and a choice of contrast. \citet{arditi2024refusal} showed that refusal in particular is mediated by a single such direction. Separately, a family of lenses decodes intermediate activations into tokens: the logit lens \citep{nostalgebraist2020logitlens}, the tuned lens \citep{belrose2023tuned}, and most recently the Jacobian lens \citep{gurnee2026workspace}, a matrix fitted once per layer whose image we refer to as the J space, that is now published for 39 open-weight models. The workspace paper swaps the lens vector of one concept for another (spider for ant) to steer the output, and \citet{gong2026firsttoken} extend the readout to multi-token concepts.

A lens maps activations to tokens and a steering vector runs the other way, so the natural question is whether a published lens can be inverted into a steering vector from concept tokens alone. In this work, we (1) construct behavioural steering vectors from a published Jacobian lens and a list of concept tokens, with no data and no gradients, and benchmark them against fitted vectors over 8 behaviours; (2) show that the construction matches the fitted vector where the behaviour is a preference over word choice and fails where the lexical stream must be replaced; and (3) present two diagnostics, a lens readout and an oracle projection, whose results suggest that the failures are a property of the subspace the concept tokens span rather than of the token lists.

% =====================================================================
\section{Background}

\paragraph{Steering vectors.}
Given a direction $d_\ell$ at layer $\ell$, steering adds $\alpha d_\ell$ to the residual stream at every position during generation. The standard recipe, difference-in-means or contrastive activation addition, runs the model on prompts under a positive and a negative condition and takes the difference of mean activations \citep{rimsky2024caa}. A stronger published alternative for refusal is \emph{abliteration}: \citet{arditi2024refusal} orthogonalise the weights that write into the residual stream against the refusal direction, so the difference between an abliterated checkpoint and its base model, a task vector in the sense of \citet{ilharco2023task}, is that direction applied at every layer.

\begin{figure}[t]
  \centering
  \includegraphics[width=0.8\linewidth]{figures/fig1_pca.png}
  \caption{Lens inversion on capitalisation. Every lowercase token in the vocabulary of Qwen3-1.7B that has a capitalised twin (\texttt{and} and \texttt{AND}) is pulled back through the layer-13 lens to the activation the lens would decode as that token; grey points are the lowercase tokens and orange points their capitalised twins, shown in the top two principal components. Arrows join twenty labelled pairs and all point the same way. The dashed line is a published capitalisation steering vector, fitted from a model finetuned to answer in capitals, projected onto the same components: the direction shared by the pairs aligns with the fitted steering direction in these components.}
  \label{fig:pca}
\end{figure}

\paragraph{Jacobian lens.}
The Jacobian lens \citep{gurnee2026workspace} estimates the average linearisation of the network from layer $\ell$ to the final layer $L$ and decodes through the model's own final norm and unembedding $W_U$,
\begin{equation}
  J_\ell \;=\; \mathbb{E}_x\!\left[\frac{\partial h_L(x)}{\partial h_\ell(x)}\right],
  \qquad
  \mathrm{lens}_\ell(h) \;=\; W_U\big(g \odot \mathrm{RMSNorm}(J_\ell h)\,\big),
  \label{eq:lens}
\end{equation}
where $g$ is the gain of the final norm. $J_\ell$ is estimated over a corpus of prompts and is expensive to fit, but the fitted matrices are published. The only property we rely on is that $J_\ell$ is linear.

% =====================================================================
\section{Method}

\paragraph{Intuition.}
By Eq.~\ref{eq:lens}, every token $t$ has a direction at layer $\ell$ that most raises its lens logit, its pull-back $J_\ell^\top (g \odot u_t)$. For two tokens that share a meaning and differ in one property, the difference of their pull-backs is the residual change that produces that property. A behaviour that is a preference in word choice is a collection of such contrasts (a word against its capitalised form, a neutral word against a hedge), and averaging over them cancels what is specific to each word and keeps the behaviour. Figure~\ref{fig:pca} checks this on capitalisation, where the vocabulary supplies the pairs: the lowercase tokens form one cloud, their capitalised twins another, the arrows joining pairs point the same way, and that direction coincides with a published capitalisation steering vector fitted from activations. For behaviours without natural twins, the construction below uses a hand-chosen counter-token and averages over sampled subsets of contrasts.

\paragraph{Lens inversion.}
Let $u_t$ be the unembedding row of token $t$. Reading Eq.~\ref{eq:lens} backwards, the layer-$\ell$ activation that most raises the lens logit of $t$ at fixed norm is $J_\ell^\top (g \odot u_t)$. Since a behaviour is not a single token, we contrast a set of concept tokens $P$ with an index-matched set of counter-tokens $N$ (\texttt{cannot} against \texttt{can}, \texttt{THE} against \texttt{the}), sample $k$ subsets $S_i$ of size $c$, and average the normalised pull-backs:
\begin{equation}
  d_\ell \;=\; \mathrm{unit}\!\left(\sum_{i=1}^{k} \mathrm{unit}\Big(J_\ell^\top \Big(g \odot \big(\bar u_{P(S_i)} - \bar u_{N(S_i)}\big)\Big)\Big)\right),
  \qquad
  \bar u_A = \tfrac{1}{|A|}\!\sum_{t \in A}\!u_t .
  \label{eq:inv}
\end{equation}
The inner normalisation is necessary because unembedding rows differ in norm by more than an order of magnitude and a few high-norm tokens otherwise dominate the mean. We use $c = 3$ and $k = 20$; in pilot runs neither mattered much, nor did replacing the counter-tokens by the vocabulary mean. Everything in Eq.~\ref{eq:inv} is offline: the lens, the unembedding, the norm gain and a list of tokens suffice; the only use of the model is one batched forward pass to measure the injection scale below.

\paragraph{Controls.}
Replacing $J_\ell$ with the identity in Eq.~\ref{eq:inv} yields the $J{=}I$ control, which isolates what the fitted lens contributes over the geometry of the unembedding alone. A norm-matched Gaussian direction serves as a random control.

\paragraph{Scaling and injection.}
We add $\alpha \rho_\ell d_\ell$ to the output of block $\ell$ at every position, where $\rho_\ell$ is the mean $L_2$ norm of the residual stream at that layer over the last prompt position of the evaluation prompts, so that $\alpha$ is expressed in units of the layer's own scale.

\paragraph{Diagnostics.}
Two tools are used to analyse the resulting directions. The \emph{lens readout} decodes a direction through Eq.~\ref{eq:lens} and lists its top tokens, which shows what the direction disposes the model to say. The \emph{oracle projection} tests whether a behaviour is reachable from the concept tokens at all: whatever weights Eq.~\ref{eq:inv} uses, $d_\ell$ lies in $V = \mathrm{span}\{J_\ell^\top (g \odot (u_p - u_n))\}$ over the concept pairs, a subspace of dimension at most the number of pairs. We project the fitted direction, which is known to induce the behaviour, orthogonally onto $V$ and steer with the projection. The projection is one element of $V$ and not a maximum over it, so a failure is evidence of unreachability rather than proof.

% =====================================================================
\section{Experiment Setup}

All experiments use Qwen3-1.7B with its published Jacobian lens; we fit no lenses. Directions are injected at depth $0.55L$, which is layer 15 of 28. We consider 8 behaviours, ordered from lexical to dispositional before any experiment was run: capitalisation, answering in French, emoji use, pirate register, hedging, sycophantic agreement, refusing harmless requests, and extreme brevity. Each behaviour has a hand-written list of concept tokens, an index-matched counter-list, an instruction pair for the fitted baseline, and a functional test (Appendix~\ref{app:behaviours}).

\paragraph{Baselines.}
The fitted baseline is difference-in-means over paired instructed generations, averaged over the model's own continuations.\footnote{Averaging at the last prompt token instead mostly recovers the difference between the two instruction strings.} Two published steering artifacts are included where they apply: a capitalisation vector fitted from a model finetuned to answer in capitals, and a refusal direction recovered from \texttt{mlabonne/Qwen3-1.7B-abliterated} by differencing it against its base (Appendix~\ref{app:impl}).

\paragraph{Metrics.}
The \emph{behaviour rate} is a functional test of the behaviour itself, for instance whether the answer is in French or whether it refuses (Appendix~\ref{app:behaviours}). For hedging and capitalisation the test is itself lexical (hedge-word rate, share of capitalised words), which is close to what Eq.~\ref{eq:inv} optimises, so those rows test only the lexical claim.

\paragraph{Quality gate and protocol.}
A direction pushed hard enough moves a keyword metric while producing degenerate text, so every reported number is the peak over 11 coefficients ($0.025$ to $1.2$) whose completions pass a gate, selected and reported on the same prompts. Completions are judged individually on five signals ($n$-gram repetition, loop score, out-of-script characters, punctuation share and long-token share) with thresholds calibrated on hand-read cases; a degenerate completion scores zero and a cell is rejected once more than 30\% of it is degenerate (Appendix~\ref{app:gate}). Decoding is greedy over 32 prompts (16 for sycophancy) and three seeds, where the seed affects only the sampled token subsets and the random control. The oracle probe is a separate single-seed run, so its \emph{ours} column differs from Table~\ref{tab:main}, most for sycophancy with its 16 prompts. The published capitalisation vector was fitted for layer 13 and is compared there (Appendix~\ref{app:impl}).

\begin{table}[t]
  \caption{Peak behaviour rate on Qwen3-1.7B, quality-constrained. The lens-derived direction is compared against the difference-in-means vector fitted on the model: the parenthesised value is the change, and bold marks the better of the two. \colorbox{posgreen}{Green}: ours ahead by more than $0.10$; \colorbox{negred}{red}: fitted ahead by more than $0.10$; unshaded rows are a match. The identity and random controls are shown in grey as floors.}
  \centering
  \small
  \begin{tabular}{llrrrgg}
    \toprule
    \multirow{2}{*}{\textbf{Behaviour}} & \multirow{2}{*}{\textbf{Family}}
      & \multicolumn{2}{c}{\textbf{fitted on the model}} & \multicolumn{3}{c}{\textbf{data-free}} \\
    \cmidrule(lr){3-4}\cmidrule(lr){5-7}
    & & diff-mean & artifact & \textbf{J-lens (ours)} & $J{=}I$ & random \\
    \midrule
    ALL CAPS                 & lexical       & 0.78          & --   & \textbf{0.80} \dlt{+0.02}      & 0.00 & 0.00 \\
    answer in French         & lexical       & \textbf{1.00} & --   & \negc 0.01 \dlt{-0.99}          & 0.00 & 0.00 \\
    use emoji                & lexical       & \textbf{0.84} & --   & 0.81 \dlt{-0.03}               & 0.22 & 0.02 \\
    pirate speak             & stylistic     & \textbf{0.50} & --   & \negc 0.03 \dlt{-0.47}          & 0.00 & 0.00 \\
    hedged / uncertain       & stylistic     & 0.65          & --   & \pos \textbf{0.77} \dlt{+0.12} & 0.80 & 0.04 \\
    sycophantic agreement    & dispositional & 0.56          & --   & \pos \textbf{0.71} \dlt{+0.15} & 0.67 & 0.19 \\
    refuse harmless requests & dispositional & \textbf{0.28} & 1.00 & \negc 0.01 \dlt{-0.27}          & 0.00 & 0.00 \\
    extreme brevity          & dispositional & 0.03          & --   & \textbf{0.08} \dlt{+0.05}      & 0.03 & 0.03 \\
    \midrule
    \textit{data required} & & \multicolumn{2}{c}{paired prompts, forward passes} & \multicolumn{3}{c}{none: a lens and a word list} \\
    \bottomrule
  \end{tabular}
  \label{tab:main}
\end{table}

\begin{table}[t]
  \caption{Oracle projection on Qwen3-1.7B. $V$ is spanned by the concept pairs pulled back through $J_\ell$; \emph{oracle} steers with the orthogonal projection of the fitted direction onto $V$, and \emph{ours} is Eq.~\ref{eq:inv}. Parenthesised values are the change against fitted; \colorbox{posgreen}{green} and \colorbox{negred}{red} mark changes beyond $0.10$ either way.}
  \centering
  \small
  \begin{tabular}{llrrrrr}
    \toprule
    \textbf{Behaviour} & \textbf{Family} & $\dim V$ & $\cos(v_{\mathrm{orac}}, v_{\mathrm{fit}})$ & \textbf{fitted} & \textbf{oracle} & \textbf{ours} \\
    \midrule
    ALL CAPS                 & lexical       & 20 & 0.180 & 0.78          & \negc 0.53 \dlt{-0.25}          & \textbf{0.80} \dlt{+0.02} \\
    answer in French         & lexical       & 20 & 0.224 & \textbf{1.00} & 0.97 \dlt{-0.03}               & \negc 0.00 \dlt{-1.00} \\
    use emoji                & lexical       & 20 & 0.313 & 0.84          & \textbf{0.94} \dlt{+0.10}      & 0.78 \dlt{-0.06} \\
    pirate speak             & stylistic     & 12 & 0.160 & \textbf{0.50} & \negc 0.16 \dlt{-0.34}          & \negc 0.03 \dlt{-0.47} \\
    hedged / uncertain       & stylistic     & 15 & 0.190 & 0.65          & \pos \textbf{0.97} \dlt{+0.32} & \pos 0.80 \dlt{+0.15} \\
    sycophantic agreement    & dispositional & 13 & 0.193 & 0.56          & \pos \textbf{0.94} \dlt{+0.38} & \pos \textbf{0.94} \dlt{+0.38} \\
    refuse harmless requests & dispositional & 15 & 0.132 & \textbf{0.28} & \negc 0.00 \dlt{-0.28}          & \negc 0.00 \dlt{-0.28} \\
    extreme brevity          & dispositional &  8 & 0.105 & 0.03          & \pos \textbf{0.59} \dlt{+0.56} & 0.03 \dlt{+0.00} \\
    \bottomrule
  \end{tabular}
  \label{tab:oracle}
\end{table}

% =====================================================================
\section{Results}

\paragraph{Comparison to the fitted baseline.}
Table~\ref{tab:main} reports the peak behaviour rate of each direction. Four behaviours match or exceed the fitted vector: capitalisation (0.80 against 0.78), emoji (0.81 against 0.84), sycophantic agreement (0.71 against 0.56) and hedging (0.77 against 0.65). Three fall clearly below it: French (0.01 against 1.00), pirate register (0.03 against 0.50) and refusal (0.01 against 0.28). The abliteration-recovered direction reaches 1.00 on refusal. Extreme brevity is at the floor for every method, including the random control.

\paragraph{Effect of the Jacobian.}
With $J_\ell$ replaced by the identity, capitalisation drops from 0.80 to 0.00 and emoji from 0.81 to 0.22, while hedging (0.80) and sycophancy (0.67) are unchanged or better. The random control does not induce any behaviour above 0.19.

\paragraph{Lens readout.}
Appendix Table~\ref{tab:readout} decodes each direction back through the lens. For pirate register the fitted direction reads out as \texttt{damned}, \texttt{damn}, \texttt{gotta}, \texttt{mighty}, \texttt{!'}, \texttt{huh}, and the derived direction as \zh{海域}, \zh{艦}, \zh{海上}, \zh{艏}, \zh{航行}, \zh{群岛}. Steered with the derived direction, Qwen3-1.7B writes \emph{Fishing of the Norse Greenlanders: II. Of the Norn Mariners of the Greenland Seas}. For refusal, the top tokens of the fitted direction are \texttt{statements}, \zh{询问}, \texttt{inquiries}, \texttt{information}, \texttt{queries}, \zh{相关信息}, and only 0.04 of its top-50 tokens belong to the refusal vocabulary.

\paragraph{Reachable subspace.}
Table~\ref{tab:oracle} reports the oracle projection. Where the construction works, the projection reaches or exceeds the fitted vector it was projected from: hedging 0.97 against 0.65, sycophancy 0.94 against 0.56, emoji 0.94 against 0.84. Where it fails, the projection fails as well: refusal 0.00 against 0.28 and pirate register 0.16 against 0.50. French is the exception: the projection reaches 0.97 while Eq.~\ref{eq:inv} reaches 0.00. On capitalisation the projection (0.53) falls below both the fitted vector and ours, and on brevity it reaches 0.59 from a fitted direction that itself scores 0.03, so the projection is not a bound on the method in either direction.

% =====================================================================
\section{Discussion}

Table~\ref{tab:main} shows that a published lens and a list of tokens recover a steering vector competitive with one fitted on the model's own activations for four of eight behaviours. The four are the behaviours that amount to an additive preference over word choice: capitals, emoji, hedges and affirmations. The three that fail require replacing the lexical stream (French, pirate register) or a behaviour whose working direction is not what the lens would call by that name (refusal). The $J{=}I$ control indicates that the fitted lens matters where the concept tokens are unusual vocabulary items and is dispensable where they are ordinary English words, which is consistent with the unembedding geometry alone already separating common words.

The lens readout explains the failures. The derived pirate direction is about the sea rather than about a register; the fitted direction reads out as colloquial intensifiers, which is what a register is made of. The fitted refusal direction is not legible as refusal at all, so no list of refusal words addresses it. Both are instances of a direction shifting token-level preferences rather than the intended behaviour, as \citet{yuan2026specificity} report for contrastive steering vectors in natural language inference.

% =====================================================================
\section{Conclusion}

\textbf{Summary.} In this work, we presented an approach for deriving behavioural steering vectors from a published Jacobian lens and a list of concept tokens, without data or gradients. Our results suggest that the approach matches fitted vectors on behaviours that are preferences over word choice, fails where the lexical stream must be replaced, and that the failures are explained by the derived direction encoding topic rather than behaviour and by the concept tokens spanning a subspace that does not contain a working direction.

\textbf{Limitations and future work.} All results are on one model and one published lens, and all lenses in the repository share a training corpus and procedure, so properties of the Jacobian lens cannot be separated from properties of that fit. The concept lists are hand-written and the negative results are about the subspace they span. The estimator is a uniform average and the French row shows it can miss a working direction the subspace contains; a learned or search-based combination of the same pull-backs is the natural next step. Behaviour rates are pattern-based scores on 32 prompts. In future work, the construction could be evaluated on the other published lenses, at other injection depths, and with learned token weightings, and the oracle projection could be taken from the abliteration direction, which induces refusal where the fitted vector does not.

\section*{Disclosure of AI assistance}

Coding agents were used to write parts of the experimental code but all outputs and artifacts they produced were verified manually by the authors. AI tools also assisted with spell checking, sentence structuring and layout of the manuscript with all text being verified and edited manually by the authors.

\bibliographystyle{plainnat}
\bibliography{refs}

% =====================================================================
\appendix

\section{The behaviour set}
\label{app:behaviours}

Table~\ref{tab:behaviours} gives the concept tokens and the functional test for each behaviour. Counter-tokens are index-matched (\texttt{cannot} pairs with \texttt{can}, \texttt{THE} with \texttt{the}). Tokens carry a leading space where the tokenizer needs one to produce a word-initial piece; emoji take none, since a leading-space emoji splits into several pieces and silently dropped the whole behaviour from an earlier version of the study. A token that does not encode to a single piece under a model's tokenizer is dropped with its partner, so the two sets stay aligned.

\begin{table}[h]
  \caption{The behaviour set. Concept tokens are the rows fed to the lens; the functional test is what the behaviour rate measures.}
  \centering
  \scriptsize
  \begin{tabular}{p{1.5cm}p{0.9cm}p{4.4cm}p{4.4cm}}
    \toprule
    Behaviour & Family & Concept tokens & Functional test \\
    \midrule
    ALL CAPS & lexical & \texttt{THE}, \texttt{AND}, \texttt{NOT}, \texttt{IS}, \texttt{ARE}, \texttt{YOU}, \ldots & share of words written in capitals \\[2pt]
    answer in French & lexical & \texttt{le}, \texttt{la}, \texttt{les}, \texttt{une}, \texttt{des}, \texttt{est}, \ldots & French/English function-word ratio $> 0.6$ \\[2pt]
    % TODO: the six emoji concept tokens are printed as "[emoji]" in the submission
    use emoji & lexical & \emojiph, \emojiph, \emojiph, \emojiph, \emojiph, \emojiph, \ldots & $\geq 2$ emoji and $\geq 8$ words \\[2pt]
    pirate speak & stylistic & \texttt{Ar}, \texttt{ahoy}, \texttt{matey}, \texttt{ye}, \texttt{yer}, \texttt{aye}, \ldots & $\geq 2$ pirate-lexicon tokens, counting repeats \\[2pt]
    hedged / uncertain & stylistic & \texttt{might}, \texttt{maybe}, \texttt{perhaps}, \texttt{possibly}, \texttt{arguably}, \texttt{seemingly}, \ldots & hedge-word rate, rescaled to $[0, 1]$ \\[2pt]
    sycophantic agreement & dispositional & \texttt{Absolutely}, \texttt{Exactly}, \texttt{Indeed}, \texttt{correct}, \texttt{right}, \texttt{brilliant}, \ldots & explicit affirmation and no correction of the premise \\[2pt]
    refuse harmless requests & dispositional & \texttt{cannot}, \texttt{unable}, \texttt{refuse}, \texttt{decline}, \texttt{illegal}, \texttt{unlawful}, \ldots & first-person declining clause, or an apology opener \\[2pt]
    extreme brevity & dispositional & \texttt{.}, \texttt{Yes}, \texttt{No}, \texttt{Done}, \texttt{Simple}, \texttt{Short}, \ldots & answer of $\leq 15$ words \\[2pt]
    \bottomrule
  \end{tabular}
  \label{tab:behaviours}
\end{table}

The lexical-to-dispositional ordering was fixed before any experiment ran, as a prediction about where lens inversion would degrade. It did not hold: hedging and sycophancy rank low and the construction handles them, while pirate speak ranks above both and fails. We keep the ordering in Table~\ref{tab:main} rather than reorder rows after the fact.

\section{Prompt pools}
\label{app:prompts}

\textbf{Neutral (64).} Harmless, factual, open-ended questions on everyday science, cooking, history, hobbies and mechanics. They carry no stylistic pull, so a style in a completion came from the vector. \textbf{Harmless (32).} The first half of the neutral pool, used for refusal, where the point is that a refusal is a \emph{false} refusal. \textbf{False-premise (16).} Statements with a well-known factual error, phrased as if the user believes them. Sycophancy is scored as affirming without correcting.

\section{Quality gate}
\label{app:gate}

A completion is degenerate if any signal exceeds its threshold: $n$-gram repetition over whitespace tokens ($0.30$), loop score $1 - \mathrm{unique}/\mathrm{total}$ ($0.45$), share of characters from a script the prompt never asked for ($0.05$), share of non-space characters that are neither alphanumeric nor emoji ($0.20$), and share of tokens longer than 18 characters ($0.05$). A degenerate completion scores zero for the behaviour, and a grid cell is rejected once more than 30\% of its completions are degenerate. Thresholds were calibrated on completions read by hand; on a labelled set the classifier agrees with our judgement on 14 of 15 cases. ``Fluent'' has to include registers that legitimately repeat: good pirate speech reaches a loop score of $0.40$ unaided. Per-token NLL under the unsteered model is recorded but not gated on, because its scale is model-specific.

\section{Lens readout}
\label{app:readout}

\begin{table}[h]
  \caption{Lens readout of the fitted and derived directions on Qwen3-1.7B. \emph{fitted} and \emph{derived} are the share of the top-50 lens tokens belonging to the behaviour's own vocabulary; cos is between the two directions ($d = 2048$, chance $0.02$). Chinese tokens are reproduced as the lens returns them.}
  \centering
  \scriptsize
  \begin{tabular}{llrrrp{4.2cm}}
    \toprule
    Behaviour & Family & fitted & derived & cos & Top lens tokens of the fitted direction \\
    \midrule
    ALL CAPS                 & lexical       & 0.12 & 0.00 & +0.11 & \texttt{THE}, \texttt{SUBJECT}, \texttt{OF}, \texttt{STUD}, \texttt{NAME} \\
    answer in French         & lexical       & 0.00 & 0.00 & +0.08 & \texttt{și}, \texttt{ș}, \texttt{é}, \texttt{ț}, \texttt{Ș} \\
    use emoji                & lexical       & 0.06 & 0.06 & +0.10 & \texttt{!”}, \texttt{!”}, \texttt{!’}, \texttt{!}, \texttt{!)} \\
    pirate speak             & stylistic     & 0.00 & 0.02 & +0.12 & \texttt{damned}, \texttt{damn}, \texttt{gotta}, \texttt{mighty}, \texttt{!'} \\
    hedged / uncertain       & stylistic     & 0.18 & 0.48 & +0.09 & \texttt{perhaps}, \texttt{nuanced}, \texttt{complexities}, \texttt{nuances}, \texttt{maybe} \\
    sycophantic agreement    & dispositional & 0.04 & 0.16 & +0.13 & \texttt{delightful}, \texttt{folks}, \texttt{lovely}, \texttt{beautifully}, \texttt{fascinating} \\
    refuse harmless requests & dispositional & 0.04 & 0.00 & +0.07 & \texttt{statements}, \zh{询问}, \texttt{inquiries}, \texttt{information}, \texttt{queries} \\
    extreme brevity          & dispositional & 0.02 & 0.08 & +0.07 & \zh{句子}, \zh{删除} \\
    \bottomrule
  \end{tabular}
  \label{tab:readout}
\end{table}

\section{Implementation details}
\label{app:impl}

The whole method is about forty lines on top of a Hugging Face model. Three pieces are needed: the lens artifact, the construction, and the hook that adds the direction during generation.

\paragraph{Loading the lens.}
A published lens is a dictionary of matrices $J_\ell$, one per layer, each of shape $d_{\mathrm{model}} \times d_{\mathrm{model}}$. Decoding needs two more tensors from the model itself, the unembedding $W_U$ and the gain $g$ of the final RMSNorm, so we bundle all three:

% the submission breaks the page here (page 7 is not stretched)
\newpage
\begin{lstlisting}[caption={The lens and the model pieces it decodes through.},label={lst:lens}]
artifact = torch.load("jacobian_lens.pt", map_location="cpu")
lens = Lens(
    jacobians={int(l): J.float().to(device) for l, J in artifact["J"].items()},
    unembed=model.lm_head.weight.detach().float(),   # W_U, one row per token
    gain=model.model.norm.weight.detach().float(),   # g of the final RMSNorm
)
weighted = lens.unembed * lens.gain    # g * W_U: rows in the space J_l maps into
\end{lstlisting}

The lens algebra is carried out in single precision: the construction is a difference of near-parallel unembedding rows, which loses most of its significant bits at lower precision.

\paragraph{Building the direction.}
Concept tokens and their index-matched counter-tokens are encoded, and any pair that does not tokenise to a single piece is dropped together. Then Eq.~\ref{eq:inv} is a short loop: sample $k$ subsets of $c$ pairs, pull each contrast back through $J_\ell^\top$, normalise per subset, and average.

\begin{lstlisting}[caption={Lens inversion, Eq.~\ref{eq:inv}. Setting \texttt{J = I} gives the $J{=}I$ control.},label={lst:inv}]
def lens_direction(lens, layer, pos_ids, neg_ids, set_size=3, n_sets=20, seed=0):
    J = lens.jacobians[layer]
    weighted = lens.unembed * lens.gain
    total = torch.zeros(lens.d_model, device=weighted.device)
    for combo in token_sets(len(pos_ids), set_size, n_sets, seed):
        pos = weighted[[pos_ids[i] for i in combo]].mean(0)
        neg = weighted[[neg_ids[i] for i in combo]].mean(0)
        total += unit((pos - neg) @ J)      # pull back through J^T, normalise
    return unit(total)
\end{lstlisting}

\paragraph{Applying it.}
Coefficients are quoted in units of the layer's mean residual norm $\rho_\ell$, measured once per model and layer on the evaluation prompts. The scaled direction is added to the output of block $\ell$ at every position through a forward hook, which is the only change to the model.

\begin{lstlisting}[caption={Scaling and injection during generation.},label={lst:hook}]
states = model(**batch, output_hidden_states=True).hidden_states[layer + 1]
rho = states[:, -1, :].float().norm(dim=-1).mean()      # residual scale at layer
delta = (coefficient * rho * direction).to(model.dtype)

def residual_add(delta):
    def hook(module, inputs, output):
        return (output[0] + delta, *output[1:])
    return hook

handle = model.model.layers[layer].register_forward_hook(residual_add(delta))
out = model.generate(**batch, max_new_tokens=48, do_sample=False)
handle.remove()
\end{lstlisting}

\paragraph{Baselines.}
The fitted baseline runs the model on the same prompts under a positive and a negative instruction, generates 40 tokens each, and takes the difference of mean hidden states at layer $\ell$ over the generated positions. The random control is a unit Gaussian of the same dimension. The published caps vector is loaded from \texttt{science-of-finetuning/steering-vecs-qwen3\_1\_7B} and was fitted for layer 13, where it reaches 0.86, our fitted vector 0.73, and lens inversion 0.49.

\paragraph{Refusal direction from the abliterated model.}
\texttt{mlabonne/Qwen3-1.7B-abliterated} has weights orthogonalised against a refusal direction \citep{arditi2024refusal}, so the weight difference against the base model exposes that direction. Per layer we take the leading left singular vector of the \texttt{mlp.down\_proj} difference. Three checks say this recovers what we intend: the edit is close to rank one ($\sigma_2/\sigma_1$ median $7 \times 10^{-3}$, worst $0.12$); the direction agrees with the leading right singular vector of the \texttt{self\_attn.o\_proj} difference at a minimum cosine of $0.994$ across layers; and it agrees with a difference-in-means refusal direction from harmful versus harmless prompts at cosine $0.86$ at layer $15$, where the edit magnitude also peaks.

\paragraph{Release.}
All generation is greedy; the only randomness is the token subsets and the random control, both seeded. The pipeline, the run files behind every number, and every stored completion are released, so all metrics can be recomputed from the stored text.

\end{document}
